\section{Methodology}\label{sec:method}

\subsection{Forecasting task and model}
Let $x_{i,t}$ be the energy (kWh) consumed by household $i$ in half-hour $t$. Each household's series is standardised
with a mean and standard deviation estimated on its own training period only. The forecaster receives the last
$L=48$ standardised readings (24\,h) and predicts the next half-hour,
$\hat{x}_{i,t+1} = f_w(x_{i,t-L+1},\ldots,x_{i,t})$. All arms share one architecture: a single-layer LSTM with 64
hidden units followed by a linear output layer ($P=\NumParams{}$ parameters, \PayloadKB{}\,kB in 32-bit floating point).
Training minimises the mean squared error with Adam (learning rate $10^{-3}$, batch 256). Errors are reported in kWh
after inverting each household's scaling. Keeping the model small and fixed is deliberate: the object of study is
the training regime, not the architecture.

\subsection{Arms}
\begin{description}
\item[Naive baselines.] Persistence ($\hat{x}_{t+1}=x_t$) and seasonal naive ($\hat{x}_{t+1}=x_{t-47}$).
\item[Local-only.] One model per household, trained on that household's windows only, with early stopping on its
own validation error. No data or weights leave the home.
\item[Centralised.] One model trained on the pooled windows of all households, with early stopping on pooled
validation error. This is the accuracy reference, and it requires the raw readings to leave the home.
\item[FedAvg.] $K$ clients each hold the windows of one or more households. In round $t$ the server samples a set
$S_t$ of clients, sends them the global weights $w_t$, each client trains for $E$ local epochs and returns $w_k$, and
the server sets $w_{t+1}=\sum_{k\in S_t'} p_k w_k$ with $p_k=n_k/\sum_{j\in S_t'} n_j$, where $n_k$ is the number of
training windows of client $k$ and $S_t'\subseteq S_t$ the clients that returned~\cite{mcmahan2017fedavg}.
\item[FedProx.] Each client minimises $F_k(w)+\frac{\mu}{2}\lVert w-w_t\rVert^2$ and stragglers' partial updates are
kept rather than dropped~\cite{li2020fedprox}.
\item[HA-FedProx.] FedProx with a coefficient $\mu_t$ set by the server each round (Section~\ref{sec:ha}).
\item[FedAvg + local fine-tuning (FT).] After federated training, each household fine-tunes the final global model
for one epoch on its own training windows. This costs no communication and is evaluated for every federated run.
\end{description}
Clients are stateless: a fresh optimiser is created at every local update, as is standard for cross-device FL.

\subsection{Clients and controlled partitions}\label{sec:partition}
Two client models are considered. In \emph{household mode} each home is a client ($K=100$), which reflects
meter-level or home-gateway deployment and has natural, uncontrolled heterogeneity. In \emph{community mode} the same
homes are grouped into $K$ clients, which reflects a building gateway, a community energy hub or an aggregator.

To dial heterogeneity with real, not invented, labels, every household is described by its normalised 24-bin mean
daily load shape and its log mean load, both computed on the training period. These features are standardised and
clustered with $k$-means ($k=5$, 10 restarts) into load-profile groups such as evening-peaking or flat, low or high consumption.
Households are then assigned to clients by one of three strategies:
(i) \emph{random}, where homes are shuffled into equal groups;
(ii) \emph{Dirichlet}$(\alpha)$~\cite{hsu2019dirichlet}, where each client draws mixture proportions
$\pi_k\sim\mathrm{Dir}(\alpha\mathbf{1})$ over profile clusters and fills its slots by sampling a cluster from
$\pi_k$, renormalised to clusters with homes left; and
(iii) \emph{cluster}, where homes are sorted by cluster and mean load and chunked, the maximally separated case.
Because the homes, and therefore the test windows, are identical under every strategy, error differences between
partitions are caused by the partition alone.

\subsection{Measured heterogeneity index}\label{sec:hindex}
Nominal labels such as ``$\alpha=0.3$'' do not say how different the resulting clients are. We therefore measure
heterogeneity directly on training data. For client $k$, let $q_k\in\Delta^{167}$ be the distribution of its total
training consumption over the 168 hours of the week. With $\mathrm{JS}$ the Jensen--Shannon divergence in
base 2~\cite{lin1991divergence} and $m_k$ the mean load of client $k$,
\begin{equation}
H_{\mathrm{JS}}=\binom{K}{2}^{-1}\sum_{k<l}\mathrm{JS}(q_k\,\Vert\,q_l),\qquad
H_{\mathrm{CV}}=\frac{\mathrm{sd}(m_k)}{\mathrm{mean}(m_k)},
\end{equation}
\begin{equation}
H=H_{\mathrm{JS}}+\frac{H_{\mathrm{CV}}}{1+H_{\mathrm{CV}}}.
\end{equation}
$H_{\mathrm{JS}}\in[0,1]$ measures differences in \emph{when} clients consume, and the bounded transform of
$H_{\mathrm{CV}}$ measures differences in \emph{how much} they consume. $H=0$ for a single client or identical
clients. The index uses no labels and no test data.

\subsection{HA-FedProx: a heterogeneity-adaptive proximal coefficient}\label{sec:ha}
The proximal coefficient of FedProx trades client drift against local progress. It is usually fixed by tuning,
yet the appropriate strength should depend on how much the client updates disagree. After aggregating round $t$,
the server holds the client updates $\Delta_k=w_k-w_t$ and their weighted mean $\bar\Delta=\sum_k p_k\Delta_k$.
It computes the \emph{update-disagreement ratio}
\begin{equation}\label{eq:rho}
\rho_t=\frac{\sum_k p_k\lVert\Delta_k-\bar\Delta\rVert^2}{\sum_k p_k\lVert\Delta_k\rVert^2}\in[0,1],
\end{equation}
which is 0 when all clients move identically and approaches 1 when their updates cancel. By the
bias--variance identity $\sum_k p_k\lVert\Delta_k\rVert^2=\lVert\bar\Delta\rVert^2+\sum_k p_k\lVert\Delta_k-\bar\Delta\rVert^2$,
$\rho_t=1-1/B_t^2$ with $B_t^2=\sum_k p_k\lVert\Delta_k\rVert^2/\lVert\bar\Delta\rVert^2$, an update-space analogue of
the $B$-local dissimilarity used in the convergence analysis of FedProx~\cite{li2020fedprox}. The next round uses
\begin{equation}
\mu_{t+1}=\mathrm{clip}\!\left(\mu_0\,\rho_t,\ \mu_{\min},\ \mu_{\max}\right),
\end{equation}
with $\mu_{\min}=10^{-4}$, $\mu_{\max}=1$ and $\mu_0$ tuned on validation data exactly as FedProx's fixed $\mu$. The
rule needs no additional client computation or messages, and it reduces to FedAvg for $\mu_0=0$. An earlier variant
normalised by $\lVert\bar\Delta\rVert^2$ alone; it is unbounded as the global model converges and
$\lVert\bar\Delta\rVert\to0$, which is why the bounded ratio~\eqref{eq:rho} is used.

\subsection{Partial participation and stragglers}
In each round a fraction $C$ of clients is sampled. Each sampled client independently becomes a straggler with
probability $p_s$ and completes only a uniform random fraction $u\sim U(0.1,1)$ of its local work. Following the
protocol of Li et al.~\cite{li2020fedprox}, FedAvg drops stragglers (they download but never upload) while FedProx
and HA-FedProx aggregate their partial updates.

\subsection{Communication accounting}\label{sec:comm}
A model message carries $P\cdot b$ bytes, with $b=4$ for 32-bit and $b=2$ for 16-bit updates. Round $t$ costs
$P b(|S_t|+|S_t'|)$: one download per selected client and one upload per returning client. Protocol framing is
excluded; the Flower backend measures it at about 1\,\% of the payload. For the centralised arm we report the volume
of the raw training readings, at 4 bytes per value plus an 8-byte timestamp, as the bytes that pooling would move.
Let $v_t$ be the validation MAE of the global model after round $t$ and $v^\star$ the centralised validation MAE.
We define
\begin{equation}
R@\varepsilon=\min\{t: v_t\le(1+\varepsilon)v^\star\},\qquad B@\varepsilon=\sum_{s\le R@\varepsilon}P b(|S_s|+|S_s'|),
\end{equation}
with $\varepsilon=10\,\%$. These convert ``how close to centralised'' into ``at what network price''.

\subsection{Metrics and statistics}
We report mean absolute error (MAE), root mean squared error (RMSE), weighted absolute percentage error
(WAPE $=100\sum|x-\hat{x}|/\sum|x|$), and MAPE restricted to half-hours with $x\ge0.05$\,kWh. Unrestricted MAPE is
unbounded at the near-zero night-time loads that are common in households. All pooled metrics are micro-averaged
over every test half-hour of every household. Per-household errors are averaged over seeds and compared with
two-sided Wilcoxon signed-rank tests~\cite{wilcoxon1945} over the \NumHouseholds{} households, with Holm
correction~\cite{holm1979} across the family of comparisons. Because $p$-values with 100 paired households can be
small for negligible effects, every test is reported with its median paired difference and the share of households
for which one arm is better.
